% SPDX-License-Identifier: LicenseRef-PULSE-NC-Research-1.0
% Generated from manuscript.md by scripts/build_pdfs.py; do not edit.
% Options for packages loaded elsewhere
\PassOptionsToPackage{unicode}{hyperref}
\PassOptionsToPackage{hyphens}{url}
\documentclass[
]{article}
\usepackage{xcolor}
\usepackage{amsmath,amssymb}
\setcounter{secnumdepth}{-\maxdimen} % remove section numbering
\usepackage{iftex}
\ifPDFTeX
  \usepackage[T1]{fontenc}
  \usepackage[utf8]{inputenc}
  \usepackage{textcomp} % provide euro and other symbols
\else % if luatex or xetex
  \usepackage{unicode-math} % this also loads fontspec
  \defaultfontfeatures{Scale=MatchLowercase}
  \defaultfontfeatures[\rmfamily]{Ligatures=TeX,Scale=1}
\fi
\usepackage{lmodern}
\ifPDFTeX\else
  % xetex/luatex font selection
\fi
% Use upquote if available, for straight quotes in verbatim environments
\IfFileExists{upquote.sty}{\usepackage{upquote}}{}
\IfFileExists{microtype.sty}{% use microtype if available
  \usepackage[]{microtype}
  \UseMicrotypeSet[protrusion]{basicmath} % disable protrusion for tt fonts
}{}
\makeatletter
\@ifundefined{KOMAClassName}{% if non-KOMA class
  \IfFileExists{parskip.sty}{%
    \usepackage{parskip}
  }{% else
    \setlength{\parindent}{0pt}
    \setlength{\parskip}{6pt plus 2pt minus 1pt}}
}{% if KOMA class
  \KOMAoptions{parskip=half}}
\makeatother
\usepackage{longtable,booktabs,array}
\usepackage{calc} % for calculating minipage widths
% Correct order of tables after \paragraph or \subparagraph
\usepackage{etoolbox}
\makeatletter
\patchcmd\longtable{\par}{\if@noskipsec\mbox{}\fi\par}{}{}
\makeatother
% Allow footnotes in longtable head/foot
\IfFileExists{footnotehyper.sty}{\usepackage{footnotehyper}}{\usepackage{footnote}}
\makesavenoteenv{longtable}
\setlength{\emergencystretch}{3em} % prevent overfull lines
\providecommand{\tightlist}{%
  \setlength{\itemsep}{0pt}\setlength{\parskip}{0pt}}
\usepackage{bookmark}
\IfFileExists{xurl.sty}{\usepackage{xurl}}{} % add URL line breaks if available
\urlstyle{same}
\hypersetup{
  pdftitle={Content-similarity ranking and topic exposure in a matched-pool recommender simulation: confirmatory effects and a like-rate boundary},
  hidelinks,
  pdfcreator={LaTeX via pandoc}}

\title{Content-similarity ranking and topic exposure in a matched-pool
recommender simulation: confirmatory effects and a like-rate boundary}
\author{}
\date{}

\begin{document}
\maketitle

\textbf{Status: independent research manuscript / preprint draft, not
peer-reviewed.} Version 0.2 (M4 synthesis, 2026-11-25, revised after the
second review pass, \texttt{research/publication/review\_pass\_2.md}).
Markdown is the canonical source; \texttt{manuscript.tex} and
\texttt{manuscript.pdf} are generated from it by
\texttt{scripts/build\_pdfs.py} (M6 release candidate, 2026-11-25).

\textbf{Author.} Mohamed Bechir Kefi (no affiliation). Authorship
approved by the owner (decision C3, 2026-11-24). An AI system was used
as a research and development assistant and is not an author (section
11).

\textbf{Data acknowledgement (required by the GroupLens licence).} This
work uses the MovieLens 1M dataset (Harper and Konstan 2015). No
endorsement by the University of Minnesota or GroupLens is stated or
implied. The data is not redistributed.

\section{Abstract}\label{abstract}

Simulation studies report that recommender feedback loops can homogenise
what users are shown. We ask how five ranking objectives (random,
popularity, content similarity, a LinUCB bandit and maximal marginal
relevance re-ranking) change exposure diversity and simulated preference
concentration when every ranker receives identical candidate pools,
users and noisy explicit feedback. Four confirmatory experiments (1,260
runs, 20 paired seeds per cell, protocol frozen before execution)
compare each ranker with random ranking. With static synthetic
preferences, content similarity lowered final-window exposure topic
entropy by 0.4000 normalised units (95\% paired t-interval {[}-0.4028,
-0.3972{]}); with MovieLens-derived preferences and catalog weights the
reduction was 0.4095 ({[}-0.4210, -0.3980{]}). Under a
drifting-preference assumption, similarity concentrated preferences by
0.4350 ({[}-0.4387, -0.4313{]}) relative to random, under which
preferences spread. A Zipf-skewed catalog enlarged the reduction instead
of removing it. Diversity-aware re-ranking recovered between 45\% and
79\% of the exposure-entropy loss across cells (supplementary Table S3),
at a small cost in relevance. All of these results hold at one
calibration of the choice curve, fitted to a like rate of 0.577 and
chosen after a pre-freeze exploratory run had shown a near-zero effect
at a like rate of about 0.02. An exploratory check made after the
confirmatory results (5 seeds) found the similarity effect near zero
when simulated users like 2\% to 5\% of what they see, and a reversed
sign under the skewed catalog. The narrowing is therefore a conditional
property of this simulation, not a general law of similarity ranking,
and nothing here describes real platforms or people.

\textbf{Keywords:} recommender systems; feedback loops; exposure
diversity; simulation; preregistration; maximal marginal relevance.

\section{1. Introduction}\label{introduction}

Recommender systems decide which items a user sees, and what the user
sees shapes the feedback the system learns from. Simulation studies have
argued that this loop can homogenise behaviour (Chaney et al.~2018),
amplify popularity bias (Mansoury et al.~2020) and degenerate user
interest (Jiang et al.~2019), and models of preference drift show
conditions under which recommendations amplify preferences (Kalimeris et
al.~2021). These results come from different simulators, rankers and
assumptions, which makes it hard to say how much of an effect belongs to
the ranking objective itself.

This study isolates the ranking objective. Five rankers receive the same
catalog, the same simulated users, the same candidate pool at every step
and the same noisy explicit feedback (common random numbers), so a
difference in exposure is attributable to the ranker within the model.
The questions, hypotheses, metrics, seeds, decision rules and one
sensitivity analysis able to restrict the interpretation were frozen
before any confirmatory run (protocol v1.0, 2026-11-24).

\textbf{Research question.} How do ranking objectives influence exposure
diversity and simulated preference reinforcement when all rankers
receive identical candidate pools?

\textbf{Hypotheses.} H1 (EXP-001): with static preferences, content
similarity yields lower final-window exposure topic entropy than random
ranking, by at least the smallest effect of interest of 0.05. H1
(EXP-002): with drifting preferences, similarity yields a larger
decrease in preference entropy than random. EXP-003 tests whether the
EXP-001 result survives a skewed catalog and a flatter choice curve;
EXP-004 repeats EXP-001 with MovieLens-derived inputs.

\textbf{Contributions.} (1) A matched-pool effect size for five standard
ranking objectives under static and drifting preference assumptions kept
strictly separate. (2) A preregistered sensitivity analysis and a second
test on real-data-derived inputs, both reported with their frozen
verdicts. (3) A disclosed, post hoc exploratory finding that the main
effect depends on how often simulated users like what they see, which
bounds the interpretation of every confirmatory result. (4) A
reproducible artifact: configs, raw run files, checksums, analysis code
committed before the data and an evidence manifest. No new phenomenon is
claimed; the direction of the primary effect was expected by
construction (protocol section 3).

\section{2. Related work}\label{related-work}

The sources below were first opened on 2026-11-24 as listed in
\texttt{research/literature\_review.md}. On 2026-11-25 (M6) the full
text of nine of the ten primary sources was opened from free, legal
copies (arXiv, author or institutional pages) and each statement made
here was checked against it; the location checked is given for each
source in \texttt{research/publication/evidence\_audit.md} (finding
F-5). The full text of Steck (2018) could not be opened legally at no
cost and is marked. This is a targeted scan, not a systematic review.

\begin{itemize}
\tightlist
\item
  \textbf{Feedback-loop homogenisation.} Chaney et al.~(2018) simulate
  recommenders trained on data confounded by their own recommendations
  and report increased homogeneity of user behaviour without a gain in
  utility (full text; sections 2 and 5.2). Mansoury et al.~(2020) report
  in offline simulation that feedback loops amplify popularity bias and
  reduce aggregate diversity (full text; section 4.3). The present study
  differs by fixing candidate pools across rankers and by measuring
  topic-level exposure rather than item-level behaviour.
\item
  \textbf{Degeneracy and drift.} Jiang et al.~(2019) formalise
  degenerate feedback loops, distinguish echo chambers from filter
  bubbles and study exploration with bandit simulations (full text;
  introduction and simulation experiments). Kalimeris et al.~(2021)
  model users drifting toward recommended content and give conditions
  for preference amplification (full text of the author PDF; sections
  2.1 and 3). EXP-002 adopts one explicit drift rule of this kind as an
  assumption; it does not test whether people drift.
\item
  \textbf{Simulation as configurable assumptions.} RecSim (Ie et
  al.~2019) and T-RECS (Lucherini et al.~2021) present recommender
  simulation as a set of configurable assumptions about users and choice
  (full texts; RecSim section 3.1, T-RECS sections 4.1 and 4.2). Our
  sensitivity analysis and the M4 exploratory check follow that view:
  they vary the assumptions most likely to drive the result.
\item
  \textbf{Rankers and metrics.} LinUCB (Li et al.~2010; full text,
  section 3.1 and Algorithm 1) is the bandit arm. Maximal marginal
  relevance (Carbonell and Goldstein 1998; full text of the two-page
  author PDF, section 2) is the diversity-aware arm. The miscalibration
  metric follows the KL definition attributed to Steck (2018), KL(p
  \textbar\textbar{} (1 - alpha) q + alpha p) with alpha 0.01. Only
  bibliographic metadata of Steck (2018) was accessed; the formula and
  the value of alpha were checked against a secondary source that
  restates them (da Silva and Jannach 2025, sections 3.1.4 and 3.1.5),
  not against Steck's own text (evidence audit, finding F-5).
\item
  \textbf{Data.} MovieLens (Harper and Konstan 2015) is cited as the
  GroupLens licence requires. The author preprint on the GroupLens site
  was opened at M6; it lists MovieLens 1M as 6,040 users, 3,706 movies
  and 1,000,209 ratings on a 1 to 5 star scale (section 3). The volume,
  issue and DOI are taken from the dataset README.
\end{itemize}

Relative to this work, PULSE is a small replication-style extension: it
reproduces the qualitative homogenisation result under tightly matched
conditions and adds effect sizes, trade-offs and a measured boundary
condition.

\section{3. Formal model}\label{formal-model}

\textbf{Items.} K = 10 topics, N = 1,000 items. Item i has primary topic
z\_i drawn from a catalog distribution pi (uniform, Zipf with exponent
1, or the MovieLens shares pi\_ML) and embedding e\_i =
normalise(onehot(z\_i) + epsilon\_i), epsilon\_i \textasciitilde{} N(0,
0.3\^{}2 I\_K).

\textbf{Users.} U = 200 users with preference vectors p\_u on the
K-simplex, either Dirichlet(alpha\_0 1\_K) with alpha\_0 in \{1.0, 0.1\}
or MovieLens-derived (section 4.3). True utility r\_ui = sum\_z p\_uz
max(e\_iz, 0) / sum\_z max(e\_iz, 0), in {[}0, 1{]}.

\textbf{Interaction loop.} For T = 200 steps, each user receives a
candidate pool of C = 50 items drawn uniformly without replacement; the
pool is identical across rankers for a given seed, user and step. The
ranker returns a slate of k = 10. The user examines position j with
probability 0.85\^{}(j-1) and likes an examined item with probability
sigmoid((r\_ui - c) / tau). Observed feedback is flipped with
probability eta.

\textbf{Preference models.} Static: p\_u(t) = p\_u(0). Drifting: after
each like of an item of topic z, p\_u \textless- normalise((1 - gamma)
p\_u + gamma onehot(z)). The two are separate assumptions and are never
pooled.

\textbf{Rankers.} (1) Random: a seeded uniform permutation of the pool.
(2) Popularity: global like count so far. (3) Content similarity: cosine
between the item embedding and the profile v\_u = (sum of liked
embeddings) - (sum of disliked embeddings), from observed feedback. (4)
LinUCB per user, ridge 1, exploration alpha\_ucb 1.0 (0.1 as a nested
level). (5) MMR re-ranking of the similarity score with lambda 0.5. All
rankers break ties with the same seeded rule.

\textbf{Metrics.} The primary metric is final-window normalised exposure
topic entropy, H\_exp(u) = -sum\_z q\_uz ln q\_uz / ln K, where q\_u is
the position-weighted topic distribution of the slates shown to u in the
last W = 50 steps; the run value is the mean over users. Lower values
mean less diverse exposure. EXP-002's primary metric is Delta H\_pref =
(H(p\_u(T)) - H(p\_u(0))) / ln K, averaged over users; negative values
mean preferences concentrated. Secondary metrics are the Gini
coefficient of item exposure counts, relevance (mean true utility of
examined items), miscalibration KL(p\_u \textbar\textbar{} 0.99 q\_u +
0.01 p\_u) in nats, user control (mean change in the exposure share of a
topic in the slate after an observed dislike of that topic, interpreted
only as arm minus random) and the observed like rate.

\section{4. Methods}\label{methods}

\subsection{4.1 Design and estimand}\label{design-and-estimand}

The experimental unit is one simulated population run (a master seed).
Users inside a run share popularity counts and the catalog, so per-user
values are averaged to one value per run before inference. All arms with
the same master seed share catalog, users and candidate pools, so
contrasts are paired by seed. The estimand is the mean over seeds of the
paired difference, treatment minus random.

\subsection{4.2 Experiments}\label{experiments}

Table 1 lists the frozen matrices. Every cell has 20 seeds (seed ranges
101001, 102001, 103001 and 104001 onward).

\textbf{Table 1. Confirmatory matrices and execution (journal
CONF-EXP-001 to CONF-EXP-004).}

\begin{longtable}[]{@{}
  >{\raggedright\arraybackslash}p{(\linewidth - 10\tabcolsep) * \real{0.1667}}
  >{\raggedright\arraybackslash}p{(\linewidth - 10\tabcolsep) * \real{0.1667}}
  >{\raggedright\arraybackslash}p{(\linewidth - 10\tabcolsep) * \real{0.1667}}
  >{\raggedright\arraybackslash}p{(\linewidth - 10\tabcolsep) * \real{0.1667}}
  >{\raggedright\arraybackslash}p{(\linewidth - 10\tabcolsep) * \real{0.1667}}
  >{\raggedright\arraybackslash}p{(\linewidth - 10\tabcolsep) * \real{0.1667}}@{}}
\toprule\noalign{}
\begin{minipage}[b]{\linewidth}\raggedright
Experiment
\end{minipage} & \begin{minipage}[b]{\linewidth}\raggedright
Question
\end{minipage} & \begin{minipage}[b]{\linewidth}\raggedright
Factors
\end{minipage} & \begin{minipage}[b]{\linewidth}\raggedright
Runs
\end{minipage} & \begin{minipage}[b]{\linewidth}\raggedright
Wall clock (s)
\end{minipage} & \begin{minipage}[b]{\linewidth}\raggedright
Failed
\end{minipage} \\
\midrule\noalign{}
\endhead
\bottomrule\noalign{}
\endlastfoot
EXP-001 & Exposure diversity, static preferences & 5 rankers x eta \{0,
0.2\} x alpha\_0 \{1.0, 0.1\}; plus bandit at alpha\_ucb 0.1 & 480 &
904.6 & 0 \\
EXP-002 & Preference reinforcement, drifting preferences & 5 rankers x
gamma \{0.01, 0.05\} x alpha\_0 \{1.0, 0.1\}, eta 0.1; plus
additive-then-clip drift check & 480 & 828.2 & 0 \\
EXP-003 & Sensitivity: catalog skew and choice sharpness & 3 rankers x
catalog \{uniform, Zipf\} x tau \{0.1, 0.3\}, eta 0.1 & 240 & 252.2 &
0 \\
EXP-004 & Second test on MovieLens-derived inputs & 3 rankers, eta 0,
tau 0.1 & 60 & 72.4 & 0 \\
\end{longtable}

\subsection{4.3 Data and calibration}\label{data-and-calibration}

Users and items in EXP-001 to EXP-003 are synthetic. MovieLens 1M (zip
SHA-256 a6898adb\ldots) enters in two ways, fixed before any rating was
read (protocol section 15). Users with odd ids form a calibration half
and users with even ids a held-out half. (a) Calibration: the target
like rate L* = 0.577143 is the share of calibration-half ratings of 4 or
more. The choice-curve midpoint c is fitted by bisection so that the
mean like probability over the synthetic world of 20 calibration seeds
equals L*, using the generators only (no ranker, no metric): c =
0.066421 for tau 0.1 and 0.005990 for tau 0.3. (b) EXP-004: preferences
are p\_u proportional to (n\_uz + 1) / pi\_ML,z, where n\_uz counts a
held-out user's ratings of 4 or more on movies of primary topic z; 200
of the 3,020 vectors are sampled per seed; the catalog follows pi\_ML;
c\_ML = 0.061067 is fitted on the calibration half. The topics are the
ten most frequent genres. \texttt{users.dat} is never read and no
row-level data is committed.

\subsection{4.4 Statistical analysis}\label{statistical-analysis}

For each contrast and metric: the mean paired difference over seeds, a
95\% paired t-interval, a percentile bootstrap interval (10,000
resamples) as a check, the standardised effect d\_z and the per-seed
values. Holm adjustment is applied within each cell and metric family.
The frozen decision rule supports H1 only if the upper bound of the
t-interval is below 0 and the point estimate is at or below -0.05 (the
smallest effect of interest). EXP-003's skew interaction uses a Welch
interval because the two catalogs differ. Nonsignificance is never read
as equivalence. The decision-rule code was committed before any
confirmatory run (decision D-013).

\subsection{4.5 Registration and
deviations}\label{registration-and-deviations}

The protocol and the four experiment protocols were frozen by the owner
on 2026-11-24 (decision D-012), with SHA-256 hashes in
\texttt{research/freeze.json}; a run guard refused any config that
differed from the freeze record. There was no deviation from the frozen
protocols during execution. Disclosed departures from an ideal record:
the reproduction-command text of three protocols is outdated (freeze
record); the experiment reports were written after the runs under an
owner decision (D-015); and the M4 exploratory check (section 5.5) was
designed after the confirmatory results were seen.

\subsection{4.6 Exploratory work before the
freeze}\label{exploratory-work-before-the-freeze}

Before the freeze, exploratory runs on disjoint seeds used the planning
midpoint c = 0.5, under which simulated users liked about 2\% of what
they saw. Under that calibration the EXP-001 primary contrast was
-0.000688 (95\% t {[}-0.002227, 0.000851{]}, 10 seeds) and the Zipf
cells of EXP-003 had the opposite sign (+0.102649 and +0.090614, 2
seeds). The owner approved recalibrating the choice curve to a stated
target like rate with these values available (decisions D-008 and
D-010). The recalibration was justified by the degenerate like rate, not
by the effect size, but it was made by people who had seen the
exploratory contrast. Section 5.5 measures what this choice implies.

\subsection{4.7 Reproducibility and
resources}\label{reproducibility-and-resources}

All runs used Python 3.12.1 and numpy 2.3.5 on an Apple M1 Pro (macOS
15.3.1), one process per experiment, at most two in parallel, code
revision f1838f3 with a clean tree. The 1,260 confirmatory runs took
about 19.5 minutes against a 4-hour budget (Table 1). Commands,
checksums and raw results are listed in section 10 and in each
experiment's \texttt{reproduction.md}.

\section{5. Results}\label{results}

Numbers in this section come from
\texttt{research/experiments/EXP-00N/\{verdict.json,\ contrasts.csv,\ summary.csv,\ tables.md\}}
and, for section 5.5, from
\texttt{results/exploratory/m4-likerate-summary.csv}. Each is mapped in
\texttt{research/publication/artifact\_manifest.json}. Intervals are
95\% paired t-intervals over 20 seeds unless stated.

\subsection{5.1 EXP-001: exposure diversity under static
preferences}\label{exp-001-exposure-diversity-under-static-preferences}

H1 was supported in the primary cell and in all three secondary cells
(Table 2, EXP-001-FIG-C01). Random ranking kept exposure close to the
pool's topic entropy (H\_exp 0.9932); similarity ranking reduced it to
0.5932.

\textbf{Table 2. EXP-001, similarity minus random in H\_exp.}

\begin{longtable}[]{@{}
  >{\raggedright\arraybackslash}p{(\linewidth - 8\tabcolsep) * \real{0.2000}}
  >{\raggedright\arraybackslash}p{(\linewidth - 8\tabcolsep) * \real{0.2000}}
  >{\raggedright\arraybackslash}p{(\linewidth - 8\tabcolsep) * \real{0.2000}}
  >{\raggedright\arraybackslash}p{(\linewidth - 8\tabcolsep) * \real{0.2000}}
  >{\raggedright\arraybackslash}p{(\linewidth - 8\tabcolsep) * \real{0.2000}}@{}}
\toprule\noalign{}
\begin{minipage}[b]{\linewidth}\raggedright
Cell
\end{minipage} & \begin{minipage}[b]{\linewidth}\raggedright
Mean
\end{minipage} & \begin{minipage}[b]{\linewidth}\raggedright
95\% t-interval
\end{minipage} & \begin{minipage}[b]{\linewidth}\raggedright
d\_z
\end{minipage} & \begin{minipage}[b]{\linewidth}\raggedright
Frozen rule
\end{minipage} \\
\midrule\noalign{}
\endhead
\bottomrule\noalign{}
\endlastfoot
eta 0.0, alpha\_0 1.0 (primary) & -0.4000 & {[}-0.4028, -0.3972{]} &
-67.47 & H1 supported \\
eta 0.0, alpha\_0 0.1 & -0.3938 & {[}-0.3963, -0.3913{]} & -74.54 & H1
supported \\
eta 0.2, alpha\_0 1.0 & -0.3959 & {[}-0.3991, -0.3927{]} & -57.30 & H1
supported \\
eta 0.2, alpha\_0 0.1 & -0.3948 & {[}-0.3974, -0.3922{]} & -71.96 & H1
supported \\
\end{longtable}

\textbf{Table 3. EXP-001 primary cell, other contrasts (mean paired
differences; Holm-adjusted p \textless{} 0.0001 unless stated).}

\begin{longtable}[]{@{}
  >{\raggedright\arraybackslash}p{(\linewidth - 10\tabcolsep) * \real{0.1667}}
  >{\raggedright\arraybackslash}p{(\linewidth - 10\tabcolsep) * \real{0.1667}}
  >{\raggedright\arraybackslash}p{(\linewidth - 10\tabcolsep) * \real{0.1667}}
  >{\raggedright\arraybackslash}p{(\linewidth - 10\tabcolsep) * \real{0.1667}}
  >{\raggedright\arraybackslash}p{(\linewidth - 10\tabcolsep) * \real{0.1667}}
  >{\raggedright\arraybackslash}p{(\linewidth - 10\tabcolsep) * \real{0.1667}}@{}}
\toprule\noalign{}
\begin{minipage}[b]{\linewidth}\raggedright
Contrast
\end{minipage} & \begin{minipage}[b]{\linewidth}\raggedright
H\_exp {[}95\% t{]}
\end{minipage} & \begin{minipage}[b]{\linewidth}\raggedright
Gini
\end{minipage} & \begin{minipage}[b]{\linewidth}\raggedright
Relevance
\end{minipage} & \begin{minipage}[b]{\linewidth}\raggedright
Miscal. KL
\end{minipage} & \begin{minipage}[b]{\linewidth}\raggedright
Like rate
\end{minipage} \\
\midrule\noalign{}
\endhead
\bottomrule\noalign{}
\endlastfoot
popularity - random & -0.0089 {[}-0.0115, -0.0064{]} & 0.7308 & 0.0003
(p 0.0022) & 0.0188 & 0.0003 (p 0.3289) \\
similarity - random & -0.4000 {[}-0.4028, -0.3972{]} & 0.1539 & 0.0603 &
0.1823 & 0.1283 \\
bandit - random & -0.2920 {[}-0.2950, -0.2891{]} & 0.1727 & 0.0536 &
0.1452 & 0.1151 \\
mmr - random & -0.1853 {[}-0.1873, -0.1834{]} & 0.0915 & 0.0503 &
-0.1874 & 0.1077 \\
mmr - similarity & 0.2147 {[}0.2125, 0.2169{]} & -0.0624 & -0.0100 &
-0.3697 & -0.0207 \\
\end{longtable}

Three observations follow. First, the popularity ranker barely changed
topic entropy (-0.0089, below the smallest effect of interest in every
cell) while it concentrated item exposure most (Gini +0.7308): topic
diversity and item concentration are different outcomes. Second,
similarity ranking had the highest relevance and like rate, so in this
model narrowing came with higher measured relevance. Third, MMR
recovered 54\% of the similarity loss (0.2147 of 0.4000; 45\% to 54\%
across the four cells, supplementary Table S3) and lowered
miscalibration below random, at a relevance cost of 0.0100. Every
non-random ranker was less responsive to dislikes than random by the
user-control metric (similarity +0.0120, MMR +0.0399), but that metric
partly measures slate volatility (limitations 13). At the nested
exploration level, LinUCB with alpha\_ucb 0.1 had a lower H\_exp than
with 1.0 in all four cells (for example 0.6803 against 0.7012 in the
primary cell; descriptive only).

\subsection{5.2 EXP-002: preference reinforcement under drifting
preferences}\label{exp-002-preference-reinforcement-under-drifting-preferences}

H1 was supported in all four cells (Table 4, EXP-002-FIG-C01). In the
primary cell preference entropy rose under random ranking (+0.0971) and
fell under similarity (-0.3379). The arms ranked as they did for
exposure: similarity, bandit (-0.2796), MMR (-0.1754), popularity
(+0.0636), random. MMR recovered 37\% of the similarity contrast (MMR
minus similarity +0.1626).

\textbf{Table 4. EXP-002, similarity minus random in Delta H\_pref, and
the drift-form check.}

\begin{longtable}[]{@{}
  >{\raggedright\arraybackslash}p{(\linewidth - 6\tabcolsep) * \real{0.2500}}
  >{\raggedright\arraybackslash}p{(\linewidth - 6\tabcolsep) * \real{0.2500}}
  >{\raggedright\arraybackslash}p{(\linewidth - 6\tabcolsep) * \real{0.2500}}
  >{\raggedright\arraybackslash}p{(\linewidth - 6\tabcolsep) * \real{0.2500}}@{}}
\toprule\noalign{}
\begin{minipage}[b]{\linewidth}\raggedright
Cell
\end{minipage} & \begin{minipage}[b]{\linewidth}\raggedright
Mean
\end{minipage} & \begin{minipage}[b]{\linewidth}\raggedright
95\% t-interval
\end{minipage} & \begin{minipage}[b]{\linewidth}\raggedright
Frozen rule
\end{minipage} \\
\midrule\noalign{}
\endhead
\bottomrule\noalign{}
\endlastfoot
gamma 0.05, alpha\_0 1.0 (primary) & -0.4350 & {[}-0.4387, -0.4313{]} &
H1 supported \\
gamma 0.01, alpha\_0 1.0 & -0.4583 & {[}-0.4612, -0.4554{]} & H1
supported \\
gamma 0.01, alpha\_0 0.1 & -0.4550 & {[}-0.4576, -0.4524{]} & H1
supported \\
gamma 0.05, alpha\_0 0.1 & -0.4336 & {[}-0.4387, -0.4285{]} & H1
supported \\
Drift-form check, additive-then-clip, alpha\_0 1.0 & -0.4363 &
{[}-0.4400, -0.4326{]} & rule not met \\
Drift-form check, additive-then-clip, alpha\_0 0.1 & -0.4351 &
{[}-0.4402, -0.4300{]} & rule not met \\
\end{longtable}

The registered drift-form check did not change the conclusion, but it is
weak: for non-negative preferences, additive-then-clip followed by
normalisation equals the multiplicative rule with rate gamma / (1 +
gamma) unless the clip binds (limitations 12). It is close to a
reparametrisation of gamma, so the result says little about other drift
forms.

\subsection{5.3 EXP-003: catalog skew and choice
sharpness}\label{exp-003-catalog-skew-and-choice-sharpness}

Neither invalidation criterion was met and no sign reversal occurred
(Table 5, EXP-003-FIG-C01). Under the Zipf catalog the random arm's
entropy fell to 0.8601, but similarity fell further (0.3540 at tau 0.1),
so the gap widened. A flatter choice curve (tau 0.3, refitted midpoint)
changed the uniform-catalog gap only slightly.

\textbf{Table 5. EXP-003, similarity minus random in H\_exp, skew
interaction and MMR recovery.}

\begin{longtable}[]{@{}
  >{\raggedright\arraybackslash}p{(\linewidth - 4\tabcolsep) * \real{0.3333}}
  >{\raggedright\arraybackslash}p{(\linewidth - 4\tabcolsep) * \real{0.3333}}
  >{\raggedright\arraybackslash}p{(\linewidth - 4\tabcolsep) * \real{0.3333}}@{}}
\toprule\noalign{}
\begin{minipage}[b]{\linewidth}\raggedright
Cell
\end{minipage} & \begin{minipage}[b]{\linewidth}\raggedright
Similarity - random {[}95\% t{]}
\end{minipage} & \begin{minipage}[b]{\linewidth}\raggedright
MMR - similarity {[}95\% t{]}
\end{minipage} \\
\midrule\noalign{}
\endhead
\bottomrule\noalign{}
\endlastfoot
tau 0.1, Zipf (primary) & -0.5060 {[}-0.5169, -0.4951{]} & 0.3669
{[}0.3597, 0.3741{]} \\
tau 0.1, uniform & -0.4001 {[}-0.4028, -0.3975{]} & 0.2150 {[}0.2126,
0.2174{]} \\
tau 0.3, uniform & -0.3861 {[}-0.3893, -0.3830{]} & 0.2306 {[}0.2274,
0.2338{]} \\
tau 0.3, Zipf & -0.5365 {[}-0.5469, -0.5260{]} & 0.4256 {[}0.4151,
0.4362{]} \\
Interaction Zipf - uniform, tau 0.1 (Welch) & -0.1059 {[}-0.1170,
-0.0948{]} & \\
Interaction Zipf - uniform, tau 0.3 (Welch) & -0.1503 {[}-0.1611,
-0.1396{]} & \\
\end{longtable}

Under skew, MMR recovered more of the loss with no detectable relevance
or like-rate cost in the tau 0.3 Zipf cell (relevance -0.0003
{[}-0.0015, 0.0008{]}), and similarity's miscalibration rose sharply
(+1.1293 and +1.4988 nats). This verdict disagrees with the pre-freeze
exploratory sweep (section 4.6), which section 5.5 examines.

\subsection{5.4 EXP-004: MovieLens-derived preferences and
catalog}\label{exp-004-movielens-derived-preferences-and-catalog}

H1 was supported (Table 6, EXP-004-FIG-C01): similarity lowered H\_exp
by 0.4095 relative to random, close to EXP-001's 0.4000. Seed-to-seed
variation was larger (paired SD 0.0246 against 0.0059), so d\_z was
smaller (-16.64). MMR stayed much closer to random (-0.0864) and
recovered 79\% of the similarity loss (0.3230 of 0.4095) at a relevance
cost of 0.0013. By the cross-experiment rule, the EXP-001 interpretation
is not restricted to the synthetic preference generator.

\textbf{Table 6. EXP-004 contrasts (Holm-adjusted p \textless{} 0.0001
unless stated).}

\begin{longtable}[]{@{}
  >{\raggedright\arraybackslash}p{(\linewidth - 6\tabcolsep) * \real{0.2500}}
  >{\raggedright\arraybackslash}p{(\linewidth - 6\tabcolsep) * \real{0.2500}}
  >{\raggedright\arraybackslash}p{(\linewidth - 6\tabcolsep) * \real{0.2500}}
  >{\raggedright\arraybackslash}p{(\linewidth - 6\tabcolsep) * \real{0.2500}}@{}}
\toprule\noalign{}
\begin{minipage}[b]{\linewidth}\raggedright
Metric
\end{minipage} & \begin{minipage}[b]{\linewidth}\raggedright
Similarity - random
\end{minipage} & \begin{minipage}[b]{\linewidth}\raggedright
MMR - random
\end{minipage} & \begin{minipage}[b]{\linewidth}\raggedright
MMR - similarity
\end{minipage} \\
\midrule\noalign{}
\endhead
\bottomrule\noalign{}
\endlastfoot
H\_exp & -0.4095 {[}-0.4210, -0.3980{]} & -0.0864 {[}-0.0911, -0.0817{]}
& 0.3230 {[}0.3134, 0.3327{]} \\
Relevance & 0.0395 {[}0.0380, 0.0410{]} & 0.0382 {[}0.0370, 0.0393{]} &
-0.0013 {[}-0.0019, -0.0007{]} (p 0.0002) \\
Miscal. KL & 0.7431 {[}0.6951, 0.7910{]} & -0.2157 {[}-0.2338,
-0.1976{]} & -0.9588 {[}-1.0002, -0.9174{]} \\
Like rate & 0.0828 {[}0.0803, 0.0853{]} & 0.0812 {[}0.0791, 0.0833{]} &
-0.0016 {[}-0.0029, -0.0003{]} (p 0.0221) \\
\end{longtable}

\subsection{5.5 Exploratory: dependence on the like rate (M4, post hoc,
5
seeds)}\label{exploratory-dependence-on-the-like-rate-m4-post-hoc-5-seeds}

\textbf{This section is exploratory. It was designed after the
confirmatory results were seen, uses 5 seeds, and never changes a frozen
verdict.} It asks whether the confirmatory effects depend on the
calibration choice disclosed in section 4.6. With the settings of
EXP-003's tau 0.1 block (static, eta 0.1, alpha\_0 1.0), the midpoint
was refitted by the frozen calibration method for target true like rates
from 0.05 to 0.85, plus the planning midpoint 0.5 (about 0.02), and
crossed with the uniform and Zipf catalogs (160 runs, seeds 960001 to
960005, journal XPL-M4-LIKERATE).

\textbf{Table 7. Similarity minus random in H\_exp by target like rate
(exploratory; 95\% t-interval over 5 seeds; FIG-M4-X01).}

\begin{longtable}[]{@{}
  >{\raggedright\arraybackslash}p{(\linewidth - 4\tabcolsep) * \real{0.3333}}
  >{\raggedright\arraybackslash}p{(\linewidth - 4\tabcolsep) * \real{0.3333}}
  >{\raggedright\arraybackslash}p{(\linewidth - 4\tabcolsep) * \real{0.3333}}@{}}
\toprule\noalign{}
\begin{minipage}[b]{\linewidth}\raggedright
True like rate (target)
\end{minipage} & \begin{minipage}[b]{\linewidth}\raggedright
Uniform catalog
\end{minipage} & \begin{minipage}[b]{\linewidth}\raggedright
Zipf catalog
\end{minipage} \\
\midrule\noalign{}
\endhead
\bottomrule\noalign{}
\endlastfoot
0.02 (planning midpoint 0.5) & -0.0032 {[}-0.0051, -0.0013{]} & +0.1045
{[}0.0925, 0.1165{]} \\
0.05 & -0.0057 {[}-0.0077, -0.0038{]} & +0.1026 {[}0.0898, 0.1154{]} \\
0.15 & -0.0384 {[}-0.0432, -0.0336{]} & +0.0724 {[}0.0572, 0.0875{]} \\
0.30 & -0.2396 {[}-0.2473, -0.2318{]} & -0.1258 {[}-0.1479,
-0.1038{]} \\
0.45 & -0.3960 {[}-0.4058, -0.3863{]} & -0.3722 {[}-0.3994,
-0.3450{]} \\
0.577 (frozen calibration) & -0.3961 {[}-0.4046, -0.3877{]} & -0.5021
{[}-0.5195, -0.4847{]} \\
0.70 & -0.3910 {[}-0.4009, -0.3811{]} & -0.6340 {[}-0.6464,
-0.6215{]} \\
0.85 & -0.3937 {[}-0.4080, -0.3795{]} & -0.7106 {[}-0.7247,
-0.6965{]} \\
\end{longtable}

At the frozen calibration the check reproduces the confirmatory values
closely (-0.3961 against -0.4001; -0.5021 against -0.5060). Away from it
the picture changes. Under the uniform catalog the similarity effect is
below the smallest effect of interest at like rates of 0.15 or less and
reaches a plateau near -0.39 from 0.45 upward. Under the Zipf catalog
the sign reverses at like rates of 0.15 or less, with similarity ranking
giving more diverse exposure than random, and the reduction grows with
the like rate above 0.30. The observed like rates of the random arm
matched the targets after allowing for the 10\% feedback flips
(\texttt{m4-likerate-summary.csv}).

\section{6. Discussion}\label{discussion}

\subsection{6.1 Expected and observed}\label{expected-and-observed}

\textbf{Table 8. Prediction, observation and frozen verdict for every
registered test.}

\begin{longtable}[]{@{}
  >{\raggedright\arraybackslash}p{(\linewidth - 6\tabcolsep) * \real{0.2500}}
  >{\raggedright\arraybackslash}p{(\linewidth - 6\tabcolsep) * \real{0.2500}}
  >{\raggedright\arraybackslash}p{(\linewidth - 6\tabcolsep) * \real{0.2500}}
  >{\raggedright\arraybackslash}p{(\linewidth - 6\tabcolsep) * \real{0.2500}}@{}}
\toprule\noalign{}
\begin{minipage}[b]{\linewidth}\raggedright
Test
\end{minipage} & \begin{minipage}[b]{\linewidth}\raggedright
Prediction (protocol)
\end{minipage} & \begin{minipage}[b]{\linewidth}\raggedright
Observation
\end{minipage} & \begin{minipage}[b]{\linewidth}\raggedright
Verdict under the frozen rule
\end{minipage} \\
\midrule\noalign{}
\endhead
\bottomrule\noalign{}
\endlastfoot
EXP-001 primary & Similarity lowers H\_exp vs random, expected by
construction & -0.4000 {[}-0.4028, -0.3972{]} & H1 supported \\
EXP-001 other cells & Same contrast in each (eta, alpha\_0) cell &
-0.3938 to -0.3959 & H1 supported in all \\
EXP-001 secondary arms & Reported with Holm, no directional rule &
Bandit -0.2920, MMR -0.1853, popularity -0.0089 & No frozen verdict;
popularity practically null \\
EXP-002 primary & Similarity concentrates preferences more than random &
-0.4350 {[}-0.4387, -0.4313{]} & H1 supported \\
EXP-002 other cells & Same contrast in each (gamma, alpha\_0) cell &
-0.4336 to -0.4583 & H1 supported in all \\
EXP-002 drift form & Dependence if sign changes or abs(D) \textless{}
0.05 & -0.4363 and -0.4351 & Not dependent (check weak, limitations
12) \\
EXP-003 skew & Gap might shrink; restrict to balanced pools if criterion
met & Gap widened: -0.5060, -0.5365 & Criterion not met; no sign
reversal \\
EXP-003 sharpness & Effect might weaken at tau 0.3 & -0.3861 (uniform),
-0.5365 (Zipf) & Criterion not met \\
EXP-003 interaction & Not directional & -0.1059 and -0.1503 (Welch) &
Reported, no rule \\
EXP-004 primary & Contrast might shrink with heterogeneous real
preferences & -0.4095 {[}-0.4210, -0.3980{]} & H1 supported; not
restricted to the synthetic generator \\
\end{longtable}

Two observations were not anticipated in the protocol: preferences
spread under random ranking in EXP-002, and skew enlarged rather than
shrank the EXP-003 gap. Both are reported as measured properties of this
model.

\subsection{6.2 Competing explanations}\label{competing-explanations}

The protocols named competing explanations in advance. Their status
after the confirmatory results and the M4 check is as follows (full
table in \texttt{research/publication/interpretation\_m4.md}).

\begin{itemize}
\tightlist
\item
  \textbf{Candidate-pool composition rather than the objective.}
  Addressed within a seed by common random numbers (pools identical
  across arms) and across catalogs by EXP-003 and EXP-004. At the frozen
  calibration the effect survives skew. The M4 check shows that the
  direction under skew depends on the like rate, so pool composition and
  calibration interact.
\item
  \textbf{Feedback noise and initial concentration.} Varied as factors
  in EXP-001 and EXP-002 without changing any verdict (range -0.3938 to
  -0.4000 in EXP-001). Under drift, the like rate of the random arm was
  0.5668 at gamma 0.05 at both alpha\_0 levels, whereas under static
  preferences it differed by alpha\_0 (0.5770 and 0.5482 at eta 0). This
  is consistent with drift overwriting most of the initial preference
  within 200 steps; that reading is a hypothesis, not tested here.
\item
  \textbf{Choice-model sharpness.} EXP-003 at tau 0.3 did not remove the
  effect. Sharpness was varied at a matched like rate, so this test does
  not address the like rate itself.
\item
  \textbf{Calibration (like rate), not named in the protocols.} This is
  the explanation the confirmatory design could not address, because one
  midpoint per tau was fixed at the freeze. The M4 check (exploratory, 5
  seeds, one setting of tau, eta and alpha\_0) suggests it matters more
  than any registered factor: under the uniform catalog the similarity
  effect stayed below the smallest effect of interest at like rates of
  0.15 or less and was clearly present from 0.30 upward. The threshold
  lies somewhere between those two levels and was not located. A
  plausible mechanism, untested here, is that when most feedback is
  negative the similarity profile is dominated by subtracted embeddings
  and steers slates away from disliked topics rather than toward liked
  ones; under a skewed catalog it is pushed away from the abundant
  topic, which raises entropy.
\item
  \textbf{Drift functional form.} Not addressed: the registered check is
  close to a reparametrisation (limitations 12).
\item
  \textbf{Implementation artefacts.} The test suite (153 tests when the
  confirmatory runs were made) covers ties, feedback updates, seed
  determinism, topic distributions and zero-feedback behaviour; the
  random arm's H\_exp stayed close to the pool topic entropy in every
  experiment (for example 0.9932 against 0.9983 in EXP-001 and 0.9088
  against 0.9139 in EXP-004); the M4 audit recomputed every contrast
  from the raw files independently of the analysis code.
\end{itemize}

\subsection{6.3 Implications}\label{implications}

Within this model, three statements are supported. At the calibrated
like rate, a ranker that personalises on topic similarity narrows topic
exposure strongly (confirmatory); the exploratory check indicates this
requires simulated users to like a substantial share of what they see.
Item popularity and topic diversity are distinct: a popularity ranker
can concentrate item exposure while leaving topic exposure almost
unchanged. Diversity-aware re-ranking recovered part of the similarity
loss with small relevance costs: in exposure entropy, 0.2147 of 0.4000
in EXP-001 and 0.3230 of 0.4095 in EXP-004; in preference concentration,
0.1626 of 0.4350 in EXP-002. The magnitudes are properties of this
simulation: a single primary topic per item, slates of 10 from pools of
50, one examination model and one calibration per experiment. They are
not estimates for any real system.

\section{7. Limitations and threats to
validity}\label{limitations-and-threats-to-validity}

The full list is \texttt{research/limitations.md} (items 1 to 18). The
most consequential are:

\begin{itemize}
\tightlist
\item
  \textbf{Simulation only.} Users and items are synthetic or constructed
  from aggregates. Nothing here measures real people, real platforms,
  attention or addiction, and no claim is made that any platform uses
  these rankers.
\item
  \textbf{Calibration dependence.} Every confirmatory result holds at
  one fitted midpoint per tau. An exploratory check shows the main
  effect near zero at low like rates and reversed in sign under skew
  (section 5.5). The target like rate itself comes from MovieLens
  ratings of movies that users chose to watch, so it is a
  selected-exposure like rate, not a like rate under random exposure
  (limitations 9, 15 to 17).
\item
  \textbf{Expected direction.} The primary direction was expected by
  construction; the informative parts are the effect sizes, trade-offs
  and boundary conditions.
\item
  \textbf{Construct validity.} H\_exp counts only an item's primary
  topic. User control mixes responsiveness and slate volatility.
  Relevance uses a latent utility that no real system observes. The
  drift-form check is weak. The miscalibration formula has been checked
  only against a secondary source, not against Steck (2018).
\item
  \textbf{Statistical.} Seed-level variance is small (paired SDs about
  0.006 to 0.025), so large d\_z values and tiny p-values follow even
  for practically null contrasts such as popularity; the smallest effect
  of interest is the relevant threshold. EXP-004's seeds resample one
  pool of 3,020 vectors, so inference is conditional on that population.
\item
  \textbf{Post hoc analysis.} Section 5.5 was designed after the results
  and has 5 seeds per cell. Its role is to bound the confirmatory
  interpretation, not to establish a new finding.
\item
  \textbf{Calibration chosen with an exploratory null in view.} The
  recalibration that produced a moderate like rate was approved after a
  pre-freeze exploratory run showed a near-zero primary contrast under
  the planning midpoint (section 4.6). The stated reason was the
  degenerate like rate (about 0.02), and the method was written before
  the fit, but a reader may reasonably treat the choice of calibration
  regime as a researcher degree of freedom. Section 5.5 makes its
  consequence explicit.
\end{itemize}

\section{8. Ethics and responsible
research}\label{ethics-and-responsible-research}

No human participants were involved and no data was collected. MovieLens
1M was used under its research licence: no redistribution, no commercial
use, acknowledgement given, demographics never read, and only aggregates
committed. Simulated agents are not participants and their behaviour is
not evidence about people. The study does not build or evaluate an
engagement-maximising system. A misuse risk is that the narrowing result
is read as advice on how to narrow exposure deliberately; the results
are reported with the diversity-aware alternative and its costs to
reduce that risk. No ethics approval was needed or claimed.

\section{9. Conclusion and future
work}\label{conclusion-and-future-work}

In a matched-pool simulation, content-similarity ranking reduced topic
exposure entropy by about 0.40 normalised units relative to random
ranking, with synthetic or MovieLens-derived preferences, and
concentrated drifting simulated preferences by a similar amount.
Diversity-aware re-ranking recovered much of the loss. These results
hold at the calibrated like rate of 0.577; an exploratory check
indicates that, under the settings it covered, the exposure effect falls
below the smallest effect of interest at low like rates and reverses in
direction under a skewed catalog. The defensible conclusion is
conditional: in this model and at this calibration, similarity ranking
narrows exposure; whether it does so at other like rates is an open,
exploratory question.

Feasible extensions, each a new experiment that would need its own
frozen protocol: a confirmatory test of the like-rate dependence with
more seeds and a registered threshold (registered as EXP-005 on
2026-11-25; not frozen, not run, no data); a structurally different
drift rule (for example drift toward exposed rather than liked topics);
multi-topic items; and the embedding noise sigma\_item, which controls
how sharply a similarity profile concentrates.

\section{10. Data, code and artifact
availability}\label{data-code-and-artifact-availability}

Repository (local; public URL: {[}to be added by the owner on
release{]}): configs in \texttt{experiments/}, raw run files in
\texttt{results/raw/EXP-00N/runs/} with SHA-256 checksums in
\texttt{research/experiments/EXP-00N/run\_manifest.json}, analysis
outputs and figures in \texttt{research/experiments/EXP-00N/},
exploratory outputs in \texttt{results/exploratory/}, and the evidence
manifest in \texttt{research/publication/artifact\_manifest.json}.
Reproduction:

\begin{verbatim}
uv sync --frozen --all-groups --all-extras
PYTHONPATH=src .venv/bin/python -m pulse.cli run --experiment EXP-00N
PYTHONPATH=src .venv/bin/python scripts/analyse_experiment.py EXP-00N
PYTHONPATH=src .venv/bin/python scripts/m4_evidence_audit.py
\end{verbatim}

An offline interactive page in \texttt{demo/} shows the same results,
copied from the committed result files by
\texttt{scripts/export\_demo\_data.py} and printed exactly as in this
manuscript (checked by \texttt{tests/test\_demo\_export.py}), together
with a labelled sandbox that is illustrative and not the research
engine.

EXP-004 needs the MovieLens files, regenerated by
\texttt{scripts/download\_movielens.py} and
\texttt{scripts/prepare\_movielens.py} and checked by SHA-256. The
MovieLens data is not redistributed. The code is source-available under
a non-commercial licence (\texttt{LICENSE}), not open source.

\section{11. Funding, competing interests and AI
assistance}\label{funding-competing-interests-and-ai-assistance}

No funding. No competing interests. An AI system (Claude, Anthropic)
acted as research and development assistant: it drafted protocols, code,
analyses, reports and this manuscript under the owner's direction. The
owner approved the protocols, the freeze and every recorded decision.
Verification performed: a test suite of 153 tests at the time of the
confirmatory runs (177 at the M6 release candidate), a run guard against
the freeze record, an independent recomputation of every contrast from
the raw files and a trace of every decimal number in the manuscript and
reports to a committed result
(\texttt{research/publication/evidence\_audit.md}). This AI review is
internal quality control, not external peer review.

\section{References}\label{references}

\begin{itemize}
\tightlist
\item
  Carbonell, J. and Goldstein, J. (1998). The use of MMR,
  diversity-based reranking for reordering documents and producing
  summaries. SIGIR 1998, pp.~335 to 336. doi:10.1145/290941.291025
  (pages from the author PDF; DOI from search results, unverified).
\item
  Chaney, A. J. B., Stewart, B. M. and Engelhardt, B. E. (2018). How
  algorithmic confounding in recommendation systems increases
  homogeneity and decreases utility. RecSys 2018. arXiv:1710.11214;
  doi:10.1145/3240323.3240370.
\item
  da Silva, D. C. and Jannach, D. (2025). Calibrated recommendations:
  survey and future directions. arXiv:2507.02643 (secondary source, used
  only to check the Steck formula).
\item
  Harper, F. M. and Konstan, J. A. (2015). The MovieLens datasets:
  history and context. ACM Transactions on Interactive Intelligent
  Systems 5(4). doi:10.1145/2827872 (author preprint opened; volume,
  issue and DOI from the GroupLens README).
\item
  Ie, E. et al.~(2019). RecSim: a configurable simulation platform for
  recommender systems. arXiv:1909.04847.
\item
  Jiang, R., Chiappa, S., Lattimore, T., Gyorgy, A. and Kohli, P.
  (2019). Degenerate feedback loops in recommender systems. AIES 2019.
  arXiv:1902.10730; doi:10.1145/3306618.3314288.
\item
  Kalimeris, D., Bhagat, S., Kalyanaraman, S. and Weinsberg, U. (2021).
  Preference amplification in recommender systems. KDD 2021.
  doi:10.1145/3447548.3467298.
\item
  Li, L., Chu, W., Langford, J. and Schapire, R. E. (2010). A
  contextual-bandit approach to personalized news article
  recommendation. WWW 2010. arXiv:1003.0146;
  doi:10.1145/1772690.1772758.
\item
  Lucherini, E., Sun, M., Winecoff, A. and Narayanan, A. (2021). T-RECS:
  a simulation tool to study the societal impact of recommender systems.
  arXiv:2107.08959.
\item
  Mansoury, M., Abdollahpouri, H., Pechenizkiy, M., Mobasher, B. and
  Burke, R. (2020). Feedback loop and bias amplification in recommender
  systems. arXiv:2007.13019.
\item
  Steck, H. (2018). Calibrated recommendations. RecSys 2018, pp.~154 to
  162. doi:10.1145/3240323.3240372 (metadata only; full text not
  opened).
\end{itemize}

\section{Supplementary appendix}\label{supplementary-appendix}

See \texttt{research/paper/supplementary.md}: the complete matrix,
configurations and hashes, extended tables, robustness results, report
links and the list of deviations and owner decisions.

\end{document}
